\documentclass[12pt,a4paper]{article}
\usepackage[margin=25mm, headheight=15pt, headsep=10pt]{geometry}
\usepackage{fontspec}
\usepackage[table]{xcolor} % Note: table option is required for rowcolors
\usepackage{titlesec}
\usepackage{tocloft}
\usepackage{fancyhdr}
\usepackage{booktabs}
\usepackage{tcolorbox}
\tcbuselibrary{skins, breakable}
\usepackage[colorlinks=true, linkcolor=emerald, urlcolor=emerald, bookmarks=true]{hyperref}

\definecolor{emerald}{RGB}{0,102,51}
\definecolor{darkred}{RGB}{139,0,0}
\definecolor{lightgray}{RGB}{245,245,245}

\usepackage[nil]{babel}
\babelprovide[import, main]{bengali}
\babelprovide[import]{arabic}
\babelfont{rm}[Renderer=HarfBuzz,Scale=1.0]{Noto Serif Bengali}
\babelfont{sf}[Renderer=HarfBuzz]{Noto Sans Bengali}
\babelfont{tt}[Renderer=HarfBuzz]{Noto Sans Bengali}
\babelfont[arabic]{rm}[Renderer=HarfBuzz,Scale=1.1]{Noto Naskh Arabic}

% Section styling
\titleformat{\section}{\Large\bfseries\color{emerald}}{\thesection.}{0.5em}{}
\titleformat{\subsection}{\large\bfseries\color{emerald!85!black}}{\thesubsection.}{0.5em}{}
\titleformat{\subsubsection}{\normalsize\bfseries\color{emerald!70!black}}{\thesubsubsection.}{0.5em}{}
\titlespacing*{\section}{0pt}{18pt}{8pt}

% Fancy Header/Footer setup
\pagestyle{fancy}
\fancyhf{} % clear all header and footer fields
\fancyhead[L]{\color{emerald}\small\textbf{\nouppercase{\leftmark}}}
\fancyfoot[L]{\scriptsize\color{black!50}{\lat{AI}}-সংকলিত, আলেম-যাচাইকৃত নয়}
\fancyfoot[C]{\color{emerald!70!black}\thepage}
\renewcommand{\headrulewidth}{0.4pt}
\renewcommand{\headrule}{\hbox to\headwidth{\color{emerald}\leaders\hrule height \headrulewidth\hfill}}
\renewcommand{\footrulewidth}{0pt}

% Custom boxes
% 1. Ayat/Hadith Box
\newtcolorbox{ayatbox}[1][]{
    enhanced, breakable, 
    colback=emerald!5, 
    colframe=emerald, 
    boxrule=0.5pt, 
    arc=3pt, 
    left=10pt, right=10pt, top=10pt, bottom=10pt,
    #1
}

% 2. Quote/Translation Box
\newtcolorbox{quotebox}[1][]{
    enhanced, breakable,
    frame hidden,
    colback=lightgray,
    borderline west={3pt}{0pt}{emerald},
    left=15pt, right=10pt, top=10pt, bottom=10pt,
    sharp corners,
    #1
}

% 3. AI-generated notice box
\newtcolorbox{warnbox}[1][]{
    enhanced, breakable,
    colback=red!4,
    colframe=darkred,
    boxrule=0.8pt,
    arc=3pt,
    left=10pt, right=10pt, top=10pt, bottom=10pt,
    #1
}

% Commands
\newcommand{\arab}[1]{\foreignlanguage{arabic}{#1}}
\newcommand{\kib}[1]{\textcolor{emerald}{\textbf{#1}}}

% Latin (English) fallback: Noto Serif Bengali has NO Latin glyphs
\newfontfamily\latfont[Renderer=HarfBuzz]{Noto Serif}
\newcommand{\lat}[1]{{\latfont #1}}

\setlength{\parskip}{5pt}
\setlength{\parindent}{0pt}

% Fix for missing bullet in Noto Serif Bengali
\renewcommand{\labelitemi}{$\bullet$}



\begin{document}
\begin{center}{\Large\bfseries\color{emerald} \lat{03}-তাউয়্যুয-বিসমিল্লাহ}\end{center}
\vspace{1em}
\section*{০৩ — তাউয়্যুয় ও বিসমিল্লাহ: \textit{\lat{A'udhu} + \lat{Bismillah}} — শরণাগতি ও আল্লাহর নামে শুরু}

\begin{warnbox}
 \textbf{গুরুত্বপূর্ণ নোটিশ:} এই কন্টেন্টটি \lat{AI} দিয়ে প্রস্তুত। কেবল সহীহ/হাসান হাদিস রাখা হয়েছে। আলেম-যাচাইকৃত নয়।
\end{warnbox}

\medskip\hrule\medskip

\subsection*{১. সূত্র কার্ড}

\begin{table}[h]\centering\small
\begin{tabular}{ll}
বিষয় & বিবরণ \\
\hline
\textbf{বিষয়} & ফাতিহা পড়ার আগে আয়াতুল কুরসির শুরুর অংশ (তাউয়্যুয়) ও বিসমিল্লাহ \\
\textbf{সূত্র} & সুনানে আবু দাউদ ১৪৫, সুনানে নাসায়ী ১৬১ (সহীহ — আলবানি) \\
\textbf{মর্যাদা} & \textbf{সুন্নাতে মুআক্কাদাহ} — ফাতিহার আগে \\
\hline
\end{tabular}
\end{table}

\medskip\hrule\medskip

\subsection*{২. মূল আরবি}

\subsubsection*{\textbf{তাউয়্যুয় (\lat{A'udhu}):}}
\begin{center}\arab{أَعُوذُ بِاللَّهِ مِنَ الشَّيْطَانِ الرَّجِيمِ}\end{center}

\subsubsection*{\textbf{বিসমিল্লাহ (\lat{Basmalah}):}}
\begin{center}\arab{بِسْمِ اللَّهِ الرَّحْمَٰنِ الرَّحِيمِ}\end{center}

\medskip\hrule\medskip

\subsection*{৩. বাংলা উচ্চারণ}

\textbf{আউযু বিল্লাহি মিনাশ শাইতানির রাজীম}
\textbf{বিসমিল্লাহির রাহমানির রাহীম}

\medskip\hrule\medskip

\subsection*{৪. শব্দ-শব্দ অর্থ}

\begin{table}[h]\centering\small
\begin{tabular}{lll}
আরবি & বাংলা & \lat{English} \\
\hline
\textbf{\arab{أَعُوذُ}} & আমি আশ্রয় চাই & \textit{\lat{I} \lat{seek} \lat{refuge}} \\
\textbf{\arab{بِاللَّهِ}} & আল্লাহর কাছে & \textit{\lat{in} \lat{Allah}} \\
\textbf{\arab{مِنَ} \arab{الشَّيْطَانِ}} & শয়তান থেকে & \textit{\lat{from} \lat{Satan}} \\
\textbf{\arab{الرَّجِيمِ}} & রাজিম (অভিশপ্ত/নিক্ষিপ্ত) & \textit{\lat{the} \lat{accursed}/\lat{outcast}} \\
\textbf{\arab{بِسْمِ}} & নামে & \textit{\lat{In} \lat{the} \lat{name}} \\
\textbf{\arab{اللَّهِ}} & আল্লাহর & \textit{\lat{of} \lat{Allah}} \\
\textbf{\arab{الرَّحْمَٰنِ}} & পরম করুণাময় & \textit{\lat{the} \lat{Most} \lat{Gracious}} \\
\textbf{\arab{الرَّحِيمِ}} & পরম দয়ালু & \textit{\lat{the} \lat{Most} \lat{Merciful}} \\
\hline
\end{tabular}
\end{table}

\medskip\hrule\medskip

\subsection*{৫. পূর্ণ বাংলা অনুবাদ}

\textbf{\lat{“}আমি অভিশপ্ত শয়তান থেকে আল্লাহর কাছে আশ্রয় চাই।\lat{”}}

\textbf{\lat{“}আল্লাহর নামে, যিনি পরম করুণাময়, পরম দয়ালু।\lat{”}}

\medskip\hrule\medskip

\subsection*{৬. সংক্ষিপ্ত ব্যাখ্যা}

\subsubsection*{\textbf{তাউয়্যুয় কী?}}
"\lat{A'udhu} (আউযু)" মানে "আমি আশ্রয় চাই"। শয়তানের কুমন্ত্রণা ও কুপ্রভাব থেকে আল্লাহর কাছে সুরক্ষা চাওয়া। এটি \textbf{\lat{shield} (ঢাল)} — নামাজে প্রবেশের আগে নিজেকে সুরক্ষিত করা।

\subsubsection*{\textbf{বিসমিল্লাহ কী?}}
"\lat{Bismillah}" — আল্লাহর নামে শুরু করা। সূরা ফাতিহা \textbf{শুরু হয়} বিসমিল্লাহ দিয়ে। এটি \textbf{\lat{intention} (নিয়ত)} ও \textbf{\lat{sincerity} (খুলুস)}-এর ঘোষণা — "আমি এই কাজটি শুধু আল্লাহর জন্য করছি।"

\subsubsection*{\textbf{কীভাবে পড়বেন?}}
\begin{itemize}
\item ধীরে, স্পষ্ট করে।
\item তাকবীরে তাহরিমা + সানা শেষে, হাত বাঁধার অবস্থায়।
\item মনে করুন: "আমি শয়তান থেকে বাঁচতে চাই, আল্লাহর নামে শুরু করছি।"
\end{itemize}
\medskip\hrule\medskip

\subsection*{৭. খুশুর জন্য মূল পয়েন্টস}

\begin{table}[h]\centering\small
\begin{tabular}{ll}
পয়েন্ট & কীভাবে প্রয়োগ করবেন \\
\hline
\textbf{শরণাগতি} & "আমি আল্লাহর সাথে আছি, শয়তান দূরে" — এই অনুভূতি রাখুন \\
\textbf{নিয়ত} & বিসমিল্লাহ পড়ার সময় মনে করুন: "শুধু আল্লাহর জন্য" \\
\textbf{ধীরতা} & তাড়াহুড়ো নয় — প্রতিটি শব্দের ওজন অনুভব করুন \\
\hline
\end{tabular}
\end{table}

\medskip\hrule\medskip

\subsection*{৮. সংশ্লিষ্ট হাদিস}

\begin{table}[h]\centering\small
\begin{tabular}{ll}
হাদিস & গ্রন্থ \\
\hline
\textbf{\lat{“}তোমাদের প্রত্যেকে যখন নামাজ পড়তে বসে, তখন সবকিছুর আগে আল্লাহর প্রশংসা, তাকবীর ও সানা পড়ুক।\lat{”}} & সুনানে আবু দাউদ ১৪৫ \\
\textbf{\lat{“}যখন তুমি মাসহাফ (কুরআন) খুলবে, তখন বিসমিল্লাহ পড়ো।\lat{”}} & সুনানে আবু দাউদ ১৪৬ \\
\hline
\end{tabular}
\end{table}

\medskip\hrule\medskip

\textit{শেষ আপডেট: আগস্ট ২০২৬}


\end{document}
